% GENERATED FILE. Edit canonical lessons, then run npm run generate:books.
% canonical-source-hash: fnv1a64:188eeb57b3fd51d4
% canonical-lessons: GU-C346-dakhlo, GU-C346-udaharan, GU-C346-tark, GU-C346-judu, GU-C346-sarkhavvu, GU-R346-review-339-346

\chapter{For Example}
\label{ch:gu-for-example}
\hlchaptermodality{346}

\begin{chapteropening}
\textbf{By the end of this chapter:} \emph{I can back up a point with an example.}

Complete the last lesson of chapter 346: I can back up a point with an example.
\end{chapteropening}

\section[\texorpdfstring{\gu{દાખલો} (dākhlo)}{દાખલો (dākhlo)}]{\gu{દાખલો} (dākhlo) — an example}

\label{lesson:GU-C346-dakhlo}



\begin{quote}
Before the new one: say the Gujarati for to list, then the Gujarati for to change one's stand.
\end{quote}

\subsection*{You'll want to know: \gu{દાખલો}}
\textbf{\gu{દાખલો}} — \emph{dākhlo} — \textquotedblleft{}an example\textquotedblright{}.

\begin{quote}
\gu{હું} \gu{તમને} \gu{એક} \gu{દાખલો} \gu{આપીશ}.
\end{quote}

\begin{quote}
\emph{hũ tamne ek dākhlo āpīś.}
\end{quote}

\begin{quote}
Let me give you an example.
\end{quote}

A masculine noun.

\begin{cousinweb}
\textbf{\gu{દાખલો}} (\emph{dākhlo}) is usually traced to Persian \emph{dākhila}, \textquotedblleft{}an entry, an admission\textquotedblright{}, a word of Arabic origin; \textquotedblleft{}an example\textquotedblright{} is a sense Gujarati gave it.
\end{cousinweb}

\begin{grammarlens}[title={What you've built}]
One more word for giving reasons and opinions.
\end{grammarlens}

\subsection*{Your turn}
\begin{itemize}
\raggedright
  \item \emph{Say it:} \emph{dākhlo}
  \item \emph{Say it:} \emph{dākhlo}, once more
  \item \emph{Say it:} \emph{hũ tamne ek dākhlo āpīś.} — Let me give you an example.
  \item \emph{Recall:} say the Gujarati for to list, then the Gujarati for to change one's stand, then say \emph{dākhlo} again
\end{itemize}

\begin{tcolorbox}[breakable,colback=teal!4,colframe=teal!35!black,title={Before you move on}]
What does \gu{દાખલો} mean? (\textquotedblleft{}An example\textquotedblright{}.) Say it once more.
\end{tcolorbox}
\section[\texorpdfstring{\gu{ઉદાહરણ} (udāharaṇ)}{ઉદાહરણ (udāharaṇ)}]{\gu{ઉદાહરણ} (udāharaṇ) — an instance}

\label{lesson:GU-C346-udaharan}



\begin{quote}
Before the new one: say the Gujarati for to change one's stand, then the Gujarati for an example.
\end{quote}

\subsection*{You'll want to know: \gu{ઉદાહરણ}}
\textbf{\gu{ઉદાહરણ}} — \emph{udāharaṇ} — \textquotedblleft{}an instance\textquotedblright{}.

\begin{quote}
\gu{ગુરુએ} \gu{પાઠ} \gu{સમજાવવા} \gu{માટે} \gu{એક} \gu{ઉદાહરણ} \gu{આપ્યું}.
\end{quote}

\begin{quote}
\emph{gurue pāṭh samjāvvā māṭe ek udāharaṇ āpyũ.}
\end{quote}

\begin{quote}
The teacher gave an instance to explain the lesson.
\end{quote}

A neuter noun.

\begin{cousinweb}
\textbf{\gu{ઉદાહરણ}} (\emph{udāharaṇ}) is a tatsama, borrowed from Sanskrit \emph{udāharaṇa}, \textquotedblleft{}relating, declaring\textquotedblright{} (Monier-Williams), and so a case cited to make a point.
\end{cousinweb}

\begin{grammarlens}[title={What you've built}]
One more word for giving reasons and opinions.
\end{grammarlens}

\subsection*{Your turn}
\begin{itemize}
\raggedright
  \item \emph{Say it:} \emph{udāharaṇ}
  \item \emph{Say it:} \emph{udāharaṇ}, once more
  \item \emph{Say it:} \emph{gurue pāṭh samjāvvā māṭe ek udāharaṇ āpyũ.} — The teacher gave an instance to explain the lesson.
  \item \emph{Recall:} say the Gujarati for to change one's stand, then the Gujarati for an example, then say \emph{udāharaṇ} again
\end{itemize}

\begin{tcolorbox}[breakable,colback=teal!4,colframe=teal!35!black,title={Before you move on}]
What does \gu{ઉદાહરણ} mean? (\textquotedblleft{}An instance\textquotedblright{}.) Say it once more.
\end{tcolorbox}
\section[\texorpdfstring{\gu{તર્ક} (tark)}{તર્ક (tark)}]{\gu{તર્ક} (tark) — reasoning}

\label{lesson:GU-C346-tark}



\begin{quote}
Before the new one: say the Gujarati for an example, then the Gujarati for an instance.
\end{quote}

\subsection*{You'll want to know: \gu{તર્ક}}
\textbf{\gu{તર્ક}} — \emph{tark} — \textquotedblleft{}reasoning\textquotedblright{}.

\begin{quote}
\gu{મીરાનો} \gu{તર્ક} \gu{સીધો} \gu{અને} \gu{સાચો} \gu{હતો}.
\end{quote}

\begin{quote}
\emph{Mīrāno tark sīdho ane sāco hato.}
\end{quote}

\begin{quote}
Mira's reasoning was straightforward and sound.
\end{quote}

A masculine noun.

\begin{cousinweb}
\textbf{\gu{તર્ક}} (\emph{tark}) is a tatsama, borrowed from Sanskrit \emph{tarka}, \textquotedblleft{}conjecture, reasoning\textquotedblright{} (Monier-Williams).
\end{cousinweb}

\begin{grammarlens}[title={What you've built}]
One more word for giving reasons and opinions.
\end{grammarlens}

\subsection*{Your turn}
\begin{itemize}
\raggedright
  \item \emph{Say it:} \emph{tark}
  \item \emph{Say it:} \emph{tark}, once more
  \item \emph{Say it:} \emph{Mīrāno tark sīdho ane sāco hato.} — Mira's reasoning was straightforward and sound.
  \item \emph{Recall:} say the Gujarati for an example, then the Gujarati for an instance, then say \emph{tark} again
\end{itemize}

\begin{tcolorbox}[breakable,colback=teal!4,colframe=teal!35!black,title={Before you move on}]
What does \gu{તર્ક} mean? (\textquotedblleft{}Reasoning\textquotedblright{}.) Say it once more.
\end{tcolorbox}
\section[\texorpdfstring{\gu{જુદું} (judũ)}{જુદું (judũ)}]{\gu{જુદું} (judũ) — different}

\label{lesson:GU-C346-judu}



\begin{quote}
Before the new one: say the Gujarati for an instance, then the Gujarati for reasoning.
\end{quote}

\subsection*{You'll want to know: \gu{જુદું}}
\textbf{\gu{જુદું}} — \emph{judũ} — \textquotedblleft{}different\textquotedblright{}.

\begin{quote}
\gu{મારો} \gu{મત} \gu{તમારા} \gu{મત} \gu{કરતાં} \gu{જુદો} \gu{છે}.
\end{quote}

\begin{quote}
\emph{māro mat tamārā mat kartā̃ judo che.}
\end{quote}

\begin{quote}
My view is different from yours.
\end{quote}

\begin{cousinweb}
\textbf{\gu{જુદું}} (\emph{judũ}) is usually traced to Persian \emph{judā}, \textquotedblleft{}separate, apart\textquotedblright{}, with the Gujarati adjective ending \emph{-ũ}.
\end{cousinweb}

\begin{grammarlens}[title={What you've built}]
One more word for giving reasons and opinions.
\end{grammarlens}

\subsection*{Your turn}
\begin{itemize}
\raggedright
  \item \emph{Say it:} \emph{judũ}
  \item \emph{Say it:} \emph{judũ}, once more
  \item \emph{Say it:} \emph{māro mat tamārā mat kartā̃ judo che.} — My view is different from yours.
  \item \emph{Recall:} say the Gujarati for an instance, then the Gujarati for reasoning, then say \emph{judũ} again
\end{itemize}

\begin{tcolorbox}[breakable,colback=teal!4,colframe=teal!35!black,title={Before you move on}]
What does \gu{જુદું} mean? (\textquotedblleft{}Different\textquotedblright{}.) Say it once more.
\end{tcolorbox}
\section[\texorpdfstring{\gu{સરખાવવું} (sarkhāvvũ)}{સરખાવવું (sarkhāvvũ)}]{\gu{સરખાવવું} (sarkhāvvũ) — to compare}

\label{lesson:GU-C346-sarkhavvu}



\begin{quote}
Before the new one: say the Gujarati for reasoning, then the Gujarati for different.
\end{quote}

\subsection*{You'll want to know: \gu{સરખાવવું}}
\textbf{\gu{સરખાવવું}} — \emph{sarkhāvvũ} — \textquotedblleft{}to compare\textquotedblright{}.

\begin{quote}
\gu{બે} \gu{દુકાનોના} \gu{ભાવ} \gu{સરખાવીને} \gu{પછી} \gu{ખરીદો}.
\end{quote}

\begin{quote}
\emph{be dukānonā bhāv sarkhāvīne pachī kharīdo.}
\end{quote}

\begin{quote}
Compare the prices at two shops and then buy.
\end{quote}

Past, with \gu{તેણે} (he or she …): \gu{તેણે} \gu{બંને} \gu{ભાવ} \gu{સરખાવ્યા}, \emph{teṇe banne bhāv sarkhāvyā}.

\begin{cousinweb}
\textbf{\gu{સરખાવવું}} (\emph{sarkhāvvũ}) is built on \emph{sarkhũ}, \textquotedblleft{}alike\textquotedblright{}, which CDIAL traces to Sanskrit \emph{sadṛkṣa}, \textquotedblleft{}like\textquotedblright{} (CDIAL 13119), with the causative ending \emph{-āvvũ}: to set things side by side.
\end{cousinweb}

\begin{grammarlens}[title={What you've built}]
One more word for giving reasons and opinions.
\end{grammarlens}

\subsection*{Your turn}
\begin{itemize}
\raggedright
  \item \emph{Say it:} \emph{sarkhāvvũ}
  \item \emph{Say it:} \emph{sarkhāvvũ}, once more
  \item \emph{Say it:} \emph{be dukānonā bhāv sarkhāvīne pachī kharīdo.} — Compare the prices at two shops and then buy.
  \item \emph{Recall:} say the Gujarati for reasoning, then the Gujarati for different, then say \emph{sarkhāvvũ} again
\end{itemize}

\begin{tcolorbox}[breakable,colback=teal!4,colframe=teal!35!black,title={Before you move on}]
What does \gu{સરખાવવું} mean? (\textquotedblleft{}To compare\textquotedblright{}.) Say it once more.
\end{tcolorbox}
\section[(dialogue)]{Looking back: from Cause and Effect to For Example}

\label{lesson:GU-R346-review-339-346}



\begin{quote}
Say the Gujarati for an effect and an influence.
\end{quote}

\subsection*{Your turn}
Say these aloud:

\begin{itemize}
\raggedright
  \item an effect, an influence, to result, a situation, a circumstance
  \item a purpose, a motive, an intention, an excuse, to accomplish
  \item a line of argument, a discussion, a point at issue, a context, to raise a point
  \item a dispute, a squabble, a disagreement, a refutation, to cut someone short
  \item backing, advocacy, against, in agreement, to lean towards
  \item a claim, proof, an assurance, uncertainty, to indicate
  \item a refusal, an explanation, a clear understanding, to list, to change one's stand
  \item an example, an instance, reasoning, different, to compare
\end{itemize}

\begin{tcolorbox}[breakable,colback=teal!4,colframe=teal!35!black,title={Before you move on}]
An effect? (\textbf{\gu{અસર}}, \emph{asar}.) An influence? (\textbf{\gu{પ્રભાવ}}, \emph{prabhāv}.) To result? (\textbf{\gu{નીપજવું}}, \emph{nīpajvũ}.) A situation? (\textbf{\gu{પરિસ્થિતિ}}, \emph{paristhiti}.) A circumstance? (\textbf{\gu{સંજોગ}}, \emph{sanjog}.) A purpose? (\textbf{\gu{ઉદ્દેશ}}, \emph{uddeś}.) A motive? (\textbf{\gu{હેતુ}}, \emph{hetu}.) An intention? (\textbf{\gu{ઇરાદો}}, \emph{irādo}.) An excuse? (\textbf{\gu{બહાનું}}, \emph{bahānũ}.) To accomplish? (\textbf{\gu{સાધવું}}, \emph{sādhvũ}.) A line of argument? (\textbf{\gu{દલીલ}}, \emph{dalīl}.) A discussion? (\textbf{\gu{ચર્ચા}}, \emph{carcā}.) A point at issue? (\textbf{\gu{મુદ્દો}}, \emph{muddo}.) A context? (\textbf{\gu{સંદર્ભ}}, \emph{sandarbh}.) To raise a point? (\textbf{\gu{ઉઠાવવું}}, \emph{uṭhāvvũ}.) A dispute? (\textbf{\gu{વિવાદ}}, \emph{vivād}.) A squabble? (\textbf{\gu{તકરાર}}, \emph{takrār}.) A disagreement? (\textbf{\gu{મતભેદ}}, \emph{matbhed}.) A refutation? (\textbf{\gu{ખંડન}}, \emph{khanḍan}.) To cut someone short? (\textbf{\gu{ટોકવું}}, \emph{ṭokvũ}.) Backing? (\textbf{\gu{ટેકો}}, \emph{ṭeko}.) Advocacy? (\textbf{\gu{હિમાયત}}, \emph{himāyat}.) Against? (\textbf{\gu{વિરુદ્ધ}}, \emph{viruddh}.) In agreement? (\textbf{\gu{સહમત}}, \emph{sahmat}.) To lean towards? (\textbf{\gu{ઢળવું}}, \emph{ḍhaḷvũ}.) A claim? (\textbf{\gu{દાવો}}, \emph{dāvo}.) Proof? (\textbf{\gu{સાબિતી}}, \emph{sābitī}.) An assurance? (\textbf{\gu{ખાતરી}}, \emph{khātrī}.) Uncertainty? (\textbf{\gu{સંદેહ}}, \emph{sandeh}.) To indicate? (\textbf{\gu{દર્શાવવું}}, \emph{darśāvvũ}.) A refusal? (\textbf{\gu{ઇનકાર}}, \emph{inkār}.) An explanation? (\textbf{\gu{ખુલાસો}}, \emph{khulāso}.) A clear understanding? (\textbf{\gu{ચોખવટ}}, \emph{cokhvaṭ}.) To list? (\textbf{\gu{ગણાવવું}}, \emph{gaṇāvvũ}.) To change one's stand? (\textbf{\gu{પલટવું}}, \emph{palaṭvũ}.) An example? (\textbf{\gu{દાખલો}}, \emph{dākhlo}.) An instance? (\textbf{\gu{ઉદાહરણ}}, \emph{udāharaṇ}.) Reasoning? (\textbf{\gu{તર્ક}}, \emph{tark}.) Different? (\textbf{\gu{જુદું}}, \emph{judũ}.) To compare? (\textbf{\gu{સરખાવવું}}, \emph{sarkhāvvũ}.)
\end{tcolorbox}
